\chapter{\texorpdfstring{\(L^p\)}{Lp} 空间}
\label{chap:lp-spaces}

可积函数在零测集上的取值不影响任何积分。因而 \(L^p\) 空间的元素不是某个
逐点函数，而是几乎处处相等函数的等价类。Hölder 不等式保证乘积可积，
Minkowski 不等式保证 \(p\) 次平均误差满足三角不等式，Riesz--Fischer
定理则说明由这一误差定义的 Cauchy 列不会离开空间。

\section{几乎处处等价类}

固定任意测度空间 \((X,\mathcal A,\mu)\)，不假设总测度有限或
\(\sigma\)-有限。对 \(0<p<\infty\)，记
\[
 \mathcal L^p(\mu)
 =\left\{f:X\to\mathbb R\text{ 可测}:
   \int_X|f|^p\,d\mu<\infty\right\}.
\]
关系
\[
 f\sim g\quad\Longleftrightarrow\quad f=g\ \text{几乎处处}
\]
是等价关系。先确认加法不会离开 \(\mathcal L^p\)：当 \(p\ge1\) 时，
凸性给出
\[
 |f+g|^p\le2^{p-1}(|f|^p+|g|^p);
\]
当 \(0<p<1\) 时，次可加性给出
\(|f+g|^p\le|f|^p+|g|^p\)。后一估计可由
\((u+v)^p-u^p\le v^p\) 得到，因为左侧作为 \(u>0\) 的函数递减。
两种情形的右侧均可积，且 \(\int|\alpha f|^p=|\alpha|^p\int|f|^p\)。
若更换几乎处处相等的代表，和与数乘只在相应零集的有限并上改变。
因此这两个运算可以定义在等价类上。

\begin{definition}
\[
 L^p(\mu)=\mathcal L^p(\mu)/{\sim}.
\]
对等价类 \([f]\)，定义
\[
 \|[f]\|_p=\left(\int_X|f|^p\,d\mu\right)^{1/p}.
\]
通常仍把等价类写作 \(f\)。
\end{definition}

若保留逐点函数而不取商，\(\|f\|_p=0\) 只推出 \(f=0\) 几乎处处，
不能推出处处为零。因此，在稍后证明三角不等式的 \(p\ge1\) 范围内，
它在逐点函数空间上只是半范数，取商后才有正定性。
\(0<p<1\) 时还存在三角不等式的问题，将另作讨论。逐点取值、某点处
连续等性质一般不能直接赋给一个 \(L^p\) 元素，除非另行选定代表。

\section{Young 不等式}

设 \(1<p<\infty\)，共轭指数 \(q\) 由
\[
 \frac1p+\frac1q=1
\]
确定。

\begin{lemma}[Young]
\label{lem:young}
对 \(a,b\ge0\)，
\[
 ab\le\frac{a^p}{p}+\frac{b^q}{q},
\]
等号当且仅当 \(a^p=b^q\)。
\end{lemma}

\begin{proof}
若 \(a=0\) 或 \(b=0\)，左侧为零，右侧非负，且等号要求两者均为零。
以下设 \(a,b>0\)。凸函数 \(\phi(t)=e^t\) 的加权 Jensen 不等式给
\[
 e^{\frac1p\log a^p+\frac1q\log b^q}
 \le\frac{a^p}{p}+\frac{b^q}{q}.
\]
左侧为 \(ab\)。严格凸性说明等号恰在
\(\log a^p=\log b^q\)，即 \(a^p=b^q\) 时成立。
\end{proof}

另一个证明是固定 \(s\ge0\)，求
\(\sup_{t\ge0}(st-t^p/p)\)。当 \(s>0\) 时，导数在
\(t=s^{1/(p-1)}\) 处由正变负，最大值为 \(s^q/q\)；
\(s=0\) 时最大值为零。于是 \(st-t^p/p\le s^q/q\)，
这就是同一个不等式。此处只需要一元极值计算。

\section{Hölder 不等式}

\begin{theorem}[Hölder]
\label{thm:holder}
若 \(1<p<\infty\)，\(1/p+1/q=1\)，\(f\in L^p\)、\(g\in L^q\)，则
\(fg\in L^1\)，且
\[
 \|fg\|_1\le\|f\|_p\|g\|_q.
\]
端点形式要求 \(f\in L^1\)、\(g\) 为本质有界可测函数，记
\(\|g\|_\infty=\operatorname*{ess\,sup}|g|<\infty\)，则
\[
 \|fg\|_1\le\|f\|_1\|g\|_\infty.
\]
\end{theorem}

\begin{proof}
若某个范数为零，相应函数几乎处处为零，乘积也几乎处处为零。以下设两个
范数均为正，并令
\[
 F=\frac{|f|}{\|f\|_p},\qquad
 G=\frac{|g|}{\|g\|_q}.
\]
逐点应用 Young 不等式并积分：
\[
 \int FG\le\frac1p\int F^p+\frac1q\int G^q=1.
\]
乘回两个范数即得结论。端点形式由
\(|g|\le\|g\|_\infty\) 几乎处处直接积分。
\end{proof}

\begin{corollary}[多重 Hölder]
\label{cor:multilinear-holder}
设 \(m\ge1\) 为整数，各 \(p_j\in[1,\infty]\)，各 \(f_j\in L^{p_j}\)，
其中 \(L^\infty\) 暂指本质有界可测函数的几乎处处等价类。
若 \(\sum_{j=1}^m1/p_j=1\)（约定 \(1/\infty=0\)），则
\[
 \left\|\prod_{j=1}^mf_j\right\|_1
 \le\prod_{j=1}^m\|f_j\|_{p_j}.
\]
\end{corollary}

\begin{proof}
先把所有 \(p_j=\infty\) 的因子用
\(|f_j|\le\|f_j\|_\infty\) 提出积分号。对余下的有限指数，若某个范数为
零，相应乘积几乎处处为零；否则令 \(F_j=|f_j|/\|f_j\|_{p_j}\)。加权算术--几何
平均不等式给出
\[
 \prod_jF_j
 =\prod_j(F_j^{p_j})^{1/p_j}
 \le\sum_j\frac{F_j^{p_j}}{p_j},
\]
因为 \(\sum_j1/p_j=1\)。积分后右侧等于 \(1\)，再乘回各范数即得结论。
若余下的唯一有限指数是 \(1\)，同一公式只有一项，仍然成立。
\end{proof}

\section{Minkowski 不等式}

\begin{theorem}[Minkowski]
\label{thm:minkowski}
若 \(1\le p<\infty\)，\(f,g\in L^p\)，则
\[
 \|f+g\|_p\le\|f\|_p+\|g\|_p.
\]
\end{theorem}

\begin{proof}
\(p=1\) 是积分三角不等式。设 \(p>1\)。先由
\[
 |f+g|^p\le2^{p-1}(|f|^p+|g|^p)
\]
确认 \(f+g\in L^p\)。令 \(q=p/(p-1)\)。若
\(\|f+g\|_p=0\) 无需证明；否则
\[
\begin{aligned}
 \|f+g\|_p^p
 &\le\int |f|\,|f+g|^{p-1}
      +\int |g|\,|f+g|^{p-1}\\
 &\le(\|f\|_p+\|g\|_p)
 \left(\int|f+g|^{(p-1)q}\right)^{1/q}\\
 &=(\|f\|_p+\|g\|_p)\|f+g\|_p^{p-1},
\end{aligned}
\]
其中第二行使用 Hölder，且 \((p-1)q=p\)。约去
\(\|f+g\|_p^{p-1}\) 即得结论。
\end{proof}

由齐次性、正定性和 Minkowski 不等式，\(\|\cdot\|_p\) 对
\(1\le p<\infty\) 是范数。

\section{\texorpdfstring{\(0<p<1\)}{0<p<1}}

当 \(0<p<1\) 时，凹性给出
\[
 |a+b|^p\le|a|^p+|b|^p,
\]
所以
\[
 d_p(f,g)=\int|f-g|^p\,d\mu
\]
在几乎处处等价类上是度量。然而
\(\|f\|_p=d_p(f,0)^{1/p}\) 一般不满足三角不等式。若 \(E,F\) 两两
不交且测度均为 \(1\)，则
\[
 \|\mathbf1_E+\mathbf1_F\|_p=2^{1/p}>2
 =\|\mathbf1_E\|_p+\|\mathbf1_F\|_p.
\]
所以这一表达式一般不是范数。不过，由凹性，
\[
 A^p+B^p\le2^{1-p}(A+B)^p\qquad(A,B\ge0).
\]
把 \(A=\|f\|_p,B=\|g\|_p\) 代入，再结合积分后的次可加性，得到
\[
 \|f+g\|_p\le2^{1/p-1}(\|f\|_p+\|g\|_p),
\]
称为拟范数。度量 \(d_p\) 仍足以研究收敛；下面证明其完备性，之后才可以
称这个空间为拟 Banach 空间。这里并不声称每种测度空间都不能另选范数：
有限个点组成的空间仍然可以用通常的有限维范数赋予同一拓扑。

设 \((f_n)\) 是 \(d_p\)-Cauchy 列，选子列使
\(d_p(f_{n_{k+1}},f_{n_k})\le2^{-k}\)。单调收敛定理给
\[
 \int\sum_k|f_{n_{k+1}}-f_{n_k}|^p\le\sum_k2^{-k}<\infty.
\]
因此这个正项级数几乎处处有限。其各项底数最终小于一，
而 \(0\le t\le1\) 时 \(t\le t^p\)，故
\(\sum_k|f_{n_{k+1}}-f_{n_k}|\) 也几乎处处收敛。
令 \(f\) 为相应逐点极限，在上述级数发散的可测零集上置零。
它是将同一零集上置零后的子列的点态极限，故可测。由
\(|a+b|^p\le|a|^p+|b|^p\) 和 Fatou 引理，
\[
 d_p(f,f_{n_k})\le\sum_{j=k}^\infty2^{-j}\longrightarrow0.
\]
特别地，
\(\int|f|^p\le\int|f-f_{n_1}|^p+\int|f_{n_1}|^p<\infty\)，
故 \(f\in L^p\)。再用 \(d_p\) 的三角不等式和原列的 Cauchy 性，
得到整个序列趋于 \(f\)，完备性得证。

\section{Riesz--Fischer 完备性定理}

\begin{theorem}[Riesz--Fischer]
\label{thm:riesz-fischer}
对 \(1\le p<\infty\)，\(L^p(\mu)\) 是 Banach 空间。
\end{theorem}

\begin{proof}
设 \((f_n)\) 是 \(L^p\)-Cauchy 列。选子列 \((f_{n_k})\)，使
\[
 \|f_{n_{k+1}}-f_{n_k}\|_p<2^{-k}.
\]
令
\[
 g_N=\sum_{k=1}^{N}|f_{n_{k+1}}-f_{n_k}|.
\]
由 Minkowski 不等式，
\[
 \|g_N\|_p\le\sum_{k=1}^{N}2^{-k}<1.
\]
\(g_N\uparrow g\)，单调收敛定理给
\[
 \int g^p=\lim_N\int g_N^p\le1.
\]
故 \(g<\infty\) 几乎处处，级数
\[
 \sum_{k=1}^{\infty}
 |f_{n_{k+1}}(x)-f_{n_k}(x)|
\]
几乎处处收敛。在可测零集 \(\{g=\infty\}\) 外定义
\(f(x)=\lim_kf_{n_k}(x)\)，在该零集上令 \(f=0\)。
这给出有限可测函数。由于几乎处处有 \(|f|\le|f_{n_1}|+g\)，
\[
 \int|f|^p\le2^{p-1}\left(\int|f_{n_1}|^p+\int g^p\right)<\infty,
\]
所以在估计范数差之前已经确认 \(f\in L^p\)。

对 \(m>k\)，
\[
 |f_{n_m}-f_{n_k}|
 \le\sum_{j=k}^{m-1}|f_{n_{j+1}}-f_{n_j}|.
\]
由 Minkowski 不等式，
\[
 \|f_{n_m}-f_{n_k}\|_p
 \le\sum_{j=k}^{m-1}2^{-j}.
\]
固定 \(k\)，令 \(m\to\infty\)，并对
\(|f_{n_m}-f_{n_k}|^p\) 使用 Fatou 引理，得到
\[
 \|f-f_{n_k}\|_p^p
 \le\liminf_{m\to\infty}\|f_{n_m}-f_{n_k}\|_p^p
 \le\left(\sum_{j=k}^{\infty}2^{-j}\right)^p.
\]
取 \(p\) 次根即知 \(\|f-f_{n_k}\|_p\to0\)。
即子列在 \(L^p\) 中收敛到 \(f\)。原列 Cauchy，故
\[
 \|f_n-f\|_p
 \le\|f_n-f_{n_k}\|_p+\|f_{n_k}-f\|_p\to0.
\]
\end{proof}

\begin{remark}
证明没有先猜出极限函数，而是从 Cauchy 列抽出“差的范数可求和”的子列，
再用几乎处处绝对收敛构造极限。这种快速子列方法也用于证明其他函数空间的
完备性。
\end{remark}

\section*{习题}
\addcontentsline{toc}{section}{习题}
\subsection*{A组　Young 与 Hölder}
\begin{enumerate}[label=\textbf{\thechapter.\arabic*.}]
\exerciseitem{lp-quotient-well-defined}
证明 \(\|[f]\|_p\) 与代表元无关，并证明函数等价类上的加法、数乘良定义。

\exerciseitem{young-calculus-proof}
不用 Jensen 不等式，考察
\(a^p/p+b^q/q-ab\) 的最小值，证明 Young 不等式及等号条件。

\exerciseitem{holder-equality}
求 Hölder 不等式在 \(1<p<\infty\) 时的等号条件。

\exerciseitem{holder-indicator}
设 \(1<p<\infty\)、\(f\in L^p(\mu)\)，且 \(E\) 可测、
\(\mu(E)<\infty\)。把 Hölder 不等式应用于 \(f\mathbf1_E\)，证明
\[
 \int_E|f|\le\mu(E)^{1-1/p}\|f\|_p.
\]

\end{enumerate}
\subsection*{B组　范数与完备性}
\begin{enumerate}[label=\textbf{\thechapter.\arabic*.},resume]
\exerciseitem{multilinear-holder-proof}
完整证明多重 Hölder 不等式，包括某些 \(p_j=\infty\) 的情形。

\exerciseitem{minkowski-integrability-gap}
补全 Minkowski 证明中先确认 \(f+g\in L^p\) 的步骤，并说明不能在此之前
直接把 \(\|f+g\|_p^{p-1}\) 当作有限数约去。

\exerciseitem{lp-parallelogram}
证明 \(L^2\) 满足平行四边形恒等式。对 \(p\ne2\)，在两点计数测度空间上
构造反例。

\exerciseitem{quasi-lp-metric}
证明 \(0<p<1\) 时 \(d_p(f,g)=\int|f-g|^p\) 是度量，并证明该度量完备。

\end{enumerate}
\subsection*{C组　级数与空间几何}
\begin{enumerate}[label=\textbf{\thechapter.\arabic*.},resume]
\exerciseitem{riesz-fischer-fatou-step}
详细证明 Riesz--Fischer 定理中
\[
 \|f-f_{n_k}\|_p\le\sum_{j=k}^\infty2^{-j}.
\]

\exerciseitem{lp-convergence-series}
设 \(1\le p<\infty\)，\(f_n\in L^p\)（\(n\ge1\)）。
若 \(\sum_n\|f_n\|_p<\infty\)，证明 \(\sum_nf_n\) 几乎处处绝对收敛，
并在 \(L^p\) 中收敛。对 \(0<p<1\)，证明改用条件
\(\sum_n\|f_n\|_p^p<\infty\) 可得到相同结论，并说明为什么不能直接沿用
\(\sum_n\|f_n\|_p<\infty\) 作为充分条件。

\optionalexerciseitem{clarkson-special-case}
设 \(2\le p<\infty\)，\(f,g\in L^p\)。证明
\[
 \|f+g\|_p^p+\|f-g\|_p^p
 \le 2^{p-1}(\|f\|_p^p+\|g\|_p^p),
\]
并证明：对每个 \(0<\varepsilon\le2\)，若
\(\|f\|_p,\|g\|_p\le1\) 且 \(\|f-g\|_p\ge\varepsilon\)，则
\[
 \left\|\frac{f+g}{2}\right\|_p
 \le\bigl(1-(\varepsilon/2)^p\bigr)^{1/p}<1.
\]
单位球中这种只依赖两点距离下界的中点估计称为一致凸性。

\projectexerciseitem{lp-completeness-template}
把 Riesz--Fischer 证明抽象成“范数可求和子列”的完备性模板，并用它重证
\(C_b(X)\) 在上确界范数下的完备性。
\end{enumerate}
